\section{Propósito y alcance}
\label{sec:proposito}

\subsection{La pregunta}

La campaña responde a una sola pregunta, y está formulada de modo que su
respuesta puede ser negativa:

\begin{quote}
\emph{Cuáles de las decisiones que definen un sistema de transferencia de
anotación basado en modelos de lenguaje y vecinos más cercanos están separadas
por evidencia sobre una ventana temporal sin fuga, y cuáles están simplemente
elegidas.}
\end{quote}

La distinción entre \emph{separada por evidencia} y \emph{elegida} no es
retórica. Está implementada: cada decisión del sistema recibe exactamente una
de cinco palabras, y la más fuerte de las cinco solo es alcanzable cuando se
declaró un umbral de comparación antes de medir y la comparación lo superó. La
sección~\ref{sec:grafo} define las cinco palabras y el orden en que se prueban.

De esto se sigue el compromiso central del diseño: la campaña se compromete por
adelantado a publicar como \emph{elegida} cualquier decisión que no alcance la
palabra fuerte, en lugar de presentarla como un hallazgo. Una campaña que solo
puede producir hallazgos no está midiendo nada.

\subsection{Lo que este documento no es}

No es un informe de resultados. Las cifras que contiene son de tres clases y
están marcadas como tales: poblaciones y costes medidos sobre el registro
actual, que fijan la escala; cotas de resolución, que fijan lo que es
observable; y valores de parámetro instanciados, que fijan la rejilla. No hay
ninguna afirmación comparativa entre configuraciones.

Tampoco es un plan de trabajo cerrado. Las decisiones que corresponden al
investigador están recogidas en la sección~\ref{sec:procedimiento} como puertas
que pueden fallar, con la entrada sobre la que fallan nombrada en cada caso.

\subsection{Por qué el diseño se escribe antes de los resultados}

Porque el defecto que esta campaña existe para corregir no es un error de
cálculo. Es estructural, y se enuncia en una frase: \textbf{un proyecto de esta
forma escribe sus decisiones en prosa y sus datos en tablas, y nada conecta las
dos}. Cuando eso ocurre, una prohibición escrita en un documento no impide la
comparación que prohíbe, un umbral declarado no puede negarse, y una fila
retirada sigue gobernando la lectura de un resultado porque la consulta que la
lee no menciona su estado.

La consecuencia práctica es que \emph{declarar} algo tiene que costar más que
escribirlo. En este diseño una declaración vale si y solo si tiene tres partes,
y no vale nada con dos:

\begin{enumerate}[leftmargin=1.5em,itemsep=0.25em]
  \item \textbf{Una fila}, en una columna tipada, direccionable por el código
        que debe obedecerla. Prosa en un fichero de texto no es una
        declaración; tampoco lo es una clave libre en un blob que ninguna
        consulta nombra.
  \item \textbf{Un guardia que pueda negarse}, es decir, algún camino que lea
        esa fila y \emph{levante una excepción} antes de que el número se
        produzca o la comparación se responda. Un guardia que registra un aviso,
        o que devuelve un valor por omisión, no es un guardia.
  \item \textbf{Dos pruebas}: una de que se niega cuando la regla se viola, y
        otra de que actúa cuando la regla se respeta. Las dos mitades son
        necesarias. Un guardia probado solo en la negativa puede ser una
        función que siempre levanta; probado solo en el camino feliz, puede ser
        una que nunca dispara.
\end{enumerate}

Todo cambio de instrumentación que propone la sección~\ref{sec:procedimiento}
se ajusta a esa plantilla, y la razón por la que este documento define cada
término es la misma: un guardia solo puede negarse a algo que esté definido.

\subsection{Audiencia y uso}

El documento está escrito para ser leído por un grupo que no ha seguido el
desarrollo. Por eso la sección~\ref{sec:definiciones} define todo el
vocabulario antes de usarlo, incluidos los términos que parecen obvios. Los
nombres técnicos se dejan en inglés, en \texttt{\small esta tipografía}, porque son los
identificadores que existen en el código y en el esquema de la base de datos.
Traducirlos crearía un segundo vocabulario que habría que mantener sincronizado
con el primero, y un conjunto cerrado escrito dos veces es la forma más común de
que las dos copias dejen de coincidir sin que nadie lo note.
